%=====================================================================
\section{Introduction}\label{sec:intro}
%=====================================================================

Cardiovascular diseases (CVDs) are the leading cause of death worldwide: the number of prevalent cases nearly doubled from 271 million in 1990 to 523 million in 2019, and CVD deaths rose from 12.1 million to 18.6 million over the same period~\cite{roth2020gbd}. The 12-lead electrocardiogram (ECG) is the most widely used cardiac test in this population, and deep neural networks now match or exceed non-specialist physicians on many ECG interpretation tasks~\cite{hannun2019,ribeiro2020,siontis2021}. Most published models, however, are trained and tested on a single database. Their reported accuracy is therefore an in-distribution estimate that says little about how the model behaves in a hospital whose patients, recording hardware, filtering and annotation habits differ from those seen in training. On the widely used PTB-XL benchmark~\cite{wagner2020,strodthoff2021}, for example, architectural refinements now differ by fractions of a point in macro area under the receiver operating characteristic curve (AUROC)~\cite{sharan2025,zhou_leadgraph2026}, whereas the much larger performance drop that occurs when a model is moved to another hospital~\cite{reyna2022,ballas2022} receives far less methodological attention.

A natural remedy is to train on several public databases at once. The PhysioNet/Computing in Cardiology (CinC) Challenges of 2020 and 2021 made this practical by releasing tens of thousands of 12-lead recordings from hospitals in China, Germany, Russia and the United States, all coded with SNOMED-CT diagnoses~\cite{perez2020,reyna2022}; PTB-XL alone contributes 21{,}837 records from 18{,}885 patients~\cite{wagner2020}. Pooling these sources exposes a problem that is easy to overlook: \emph{each database annotates a different subset of diagnoses}. The CPSC2018 database~\cite{liu_cpsc2018}, for instance, was labeled for nine rhythm, conduction and ST-segment classes, whereas other databases code findings such as T-wave abnormalities, axis deviation or low QRS voltage that CPSC2018 never mentions. In the four databases studied here, only 18 of the 43 diagnoses in our vocabulary are annotated by all four databases, 27.9\% of the positive labels in the pooled data belong to diagnoses that at least one database never codes, and 16.7\% of all record--diagnosis entries are unknown rather than negative (Fig.~\ref{fig:annot}). When databases are concatenated, a diagnosis that a hospital never codes is silently interpreted as a confirmed negative. The model is thereby trained to deny findings that may well be present, and the error is systematic: it is correlated with the source, and therefore with the very acquisition characteristics that a cross-hospital model needs to ignore. The organizers of the 2021 Challenge identified heterogeneous labels and populations as a central open issue~\cite{reyna2022}.

Four lines of work address parts of this problem, and each stops short of the setting we study. First, the architectural line builds ever stronger classifiers evaluated in-distribution~\cite{zhao2020cinc,nejedly2021,sharan2025,zhou_leadgraph2026,sethi2025proto}; it assumes a complete label matrix and does not test transfer to a database that contributed no training data. Second, partial-label methods treat unannotated entries as unknown rather than negative~\cite{durand2019,benbaruch2022}; for ECG, Li et al.~\cite{li2024} recovered missing super-classes with a category mapping and masked predictions for dataset-related missing labels. These methods operate inside a fixed label set, so a diagnosis that no training source annotates cannot be scored at all, and they are evaluated on test splits drawn from the training sources. Third, ECG--language models align signals with free-text reports and can score diagnoses described in natural language, with zero-shot transfer across databases~\cite{merl2024,suprime2025,melp2025}. They require hundreds of thousands of report-paired recordings and substantial pretraining compute, and they do not model which diagnoses each source actually annotates; because most public 12-lead databases carry categorical codes rather than reports, their supervision cannot be used by report-based pretraining. Finally, domain-generalization (DG) techniques~\cite{ganin2016,sun2016coral,zhou2021mixstyle,sagawa2020} have been applied to multi-source ECG~\cite{ballas2022}, but they assume that every source shares the same, fully observed label space.

This paper treats the three issues jointly and asks three research questions. (RQ1) Can a model trained on several heterogeneously annotated databases learn correctly from labels that some sources never annotate? (RQ2) Does such a model generalize to a database that contributed no training data? (RQ3) Can it score diagnoses for which no training label exists, using only a clinical description? Unlike report-pretrained ECG--language models, we obtain open-vocabulary scoring from the categorical labels that public databases already carry, by modeling which labels each source observes. Every diagnosis is represented by a short, guideline-based description, encoded once by a frozen biomedical text encoder. A one-dimensional residual encoder produces a sequence of ECG tokens that is read by two heads: a conventional linear head with one weight vector per diagnosis, and a text-anchored head in which label queries derived from the description embeddings attend over the tokens and are scored by cosine similarity with their anchors, so that this head has no diagnosis-specific weights. Diagnoses that the training sources annotate are scored by the mean of the two heads' probabilities, and diagnoses that they do not annotate by the text head alone. Both heads are trained only on the labels each source annotates, after completion of positives with a small set of clinical implication rules, and an exponential moving average of the weights is used for model selection and inference. Every compared method shares the encoder, augmentation, weight averaging and model-selection rule.

The contributions are as follows.
\begin{itemize}
\item We formalize multi-source ECG learning with source-specific label vocabularies, in which an unannotated entry is unknown rather than negative, show what a naively pooled model estimates, and quantify in a controlled experiment how treating unannotated diagnoses as negatives affects performance on retained and removed diagnoses.
\item We propose a hybrid classifier that fuses a per-diagnosis linear head with a text-anchored label-query head on a shared encoder, trained with a masked asymmetric loss and implication completion. It needs no report-paired pretraining, trains in about twelve minutes on a single cloud GPU, and can score a new diagnosis from its description alone.
\item We introduce a leave-one-database-out protocol over four public databases from three countries (87{,}661 recordings) that separates diagnoses seen during training from those annotated only in the held-out database, together with a controlled annotation-removal experiment on PTB-XL that isolates the effect of label handling from that of distribution shift.
\item We compare the method against seven label-handling and DG baselines (naive pooling, masked loss, Li-style masking with category mapping, and CORAL, DANN, GroupDRO and MixStyle on top of the latter) trained under identical conditions, with three seeds per method, four ablations of the hybrid design and multiplicity-corrected tests at the level of runs and of test records. The method matches the strongest baseline on seen diagnoses, is significantly better than six of the seven baselines, and is the only trained method that scores diagnoses without training labels above chance (macro AUROC 67.0 against 51--53). We further evaluate zero-shot scoring systematically on 120 withheld diagnosis--database pairs against a report-pretrained ECG--language model (MERL), identify when description-based transfer fails, and assess calibration and the sensitivity to the annotation threshold. We also report the development history of the method, including a text-only design that underperformed the baselines and one deviation from our pre-specified selection rule. Code, descriptions, split seeds and all run logs are released.
\end{itemize}

The remainder of the paper is organized as follows. Section~\ref{sec:related} reviews related work and positions our contribution. Section~\ref{sec:methods} formalizes the problem and describes the model and objective. Section~\ref{sec:setup} details the data, protocols, baselines, implementation and statistics. Section~\ref{sec:results} reports the results, Section~\ref{sec:discussion} discusses clinical implications and limitations, and Section~\ref{sec:conclusion} concludes.
